\begin{enumerate}
\item La EDO se verifica donde la tangente está definida, pero no en $\mathbb{R}$. Los intervalos máximos son $(-\pi/2+k\pi-C,\,\pi/2+k\pi-C)$, $k\in\mathbb{Z}$.
\item Los miembros son soluciones, pero sólo admiten $a\geq0$. Para $a>0$ hay dos valores $C=\pm\sqrt{a/2}$; para $a=0$, sólo $C=0$. La familia completa sin redundancia es $y_A=A(2+\sin x)$, $A\in\mathbb{R}$, con $A=a/2$.
\item El dominio es $\mathbb{R}$ si $C\geq0$ y $(\ln(-C),\infty)$ si $C<0$. Sólo $C=1$ produce una solución de la EDO indicada.
\item Todos los miembros son soluciones en $\mathbb{R}$, pero también lo es $z(x)=xe^{-x^2}$. Esta función no pertenece a la familia en ningún intervalo abierto no vacío, pues $z(x)/e^{-x^2}=x$ no es constante.
\item La parametrización dada recorre sólo $(0,\infty)$ y verifica la EDO. Para $x<0$ sirve $x=-e^t$, $y=C+e^{-t}$, $t\in\mathbb{R}$. En ambos casos la gráfica es $y=C-1/x$ en la componente correspondiente.
\end{enumerate}
